\documentclass[11pt]{article}
\usepackage[a4paper,margin=1in]{geometry}
\usepackage{fontspec}
\usepackage[boldfont=false]{xeCJK}
\setCJKmainfont{FandolSong-Regular.otf}
\usepackage{amsmath}
\usepackage{amssymb}
\usepackage{array}
\usepackage{booktabs}
\usepackage{graphicx}
\usepackage{adjustbox}
\usepackage{enumitem}
\setlist[itemize]{topsep=2pt,partopsep=0pt,parsep=0pt,itemsep=2pt,leftmargin=1.4em}
\setlist[enumerate]{topsep=2pt,partopsep=0pt,parsep=0pt,itemsep=2pt,leftmargin=1.6em}
\usepackage{parskip}
\usepackage{fvextra}
\fvset{breaklines=true,breakanywhere=true,fontsize=\small}
\setlength{\parindent}{0pt}

\begin{document}

\begin{center}
  {\LARGE\bfseries Carboxylic acids and derivatives}\\[0.35em]
  {\large A Level Chemistry}
\end{center}
\vspace{0.6em}

\section*{Carboxylic acids}

\subsection*{Making and reacting}

\begin{itemize}
\item an \textbf{alkylbenzene} 烷基苯 (such as methylbenzene) is oxidised by hot alkaline $\text{KMnO}_4$, then dilute acid, to give \textbf{benzoic acid} 苯甲酸. The whole side-chain becomes a $\text{–COOH}$ group.

\item a carboxylic acid reacts with $\text{PCl}_3$ and heat, $\text{PCl}_5$, or $\text{SOCl}_2$ to form an \textbf{acyl chloride} 酰氯.

\end{itemize}

\subsection*{Acids that can be oxidised further}

Two \textbf{carboxylic acids} 羧酸 are special because they can still be oxidised:

\begin{itemize}
\item methanoic acid ($\text{HCOOH}$) is oxidised by Fehling's or Tollens' reagent, or acidified $\text{KMnO}_4$ / $\text{K}_2\text{Cr}_2\text{O}_7$, to carbon dioxide and water.

\item ethanedioic acid ($\text{HOOCCOOH}$) is oxidised by warm acidified $\text{KMnO}_4$ to carbon dioxide.

\end{itemize}

\subsection*{Relative acidities}

The \textbf{acidity} 酸性 order is:

\[
\text{alcohol} < \text{phenol} < \text{carboxylic acid}
\]

A carboxylic acid is the strongest because, when it loses $\text{H}^+$, the negative charge is spread over \textbf{two} oxygen atoms, making the ion very stable. In a \textbf{phenol} 苯酚 the charge spreads only into the ring, and in an \textbf{alcohol} 醇 (giving $\text{RO}^-$) it is not spread at all.

Chlorine atoms make a carboxylic acid \textbf{more} acidic. Chlorine is \textbf{electron-withdrawing} 吸电子: through the \textbf{inductive effect} 诱导效应 it pulls electron density away, helping to spread the negative charge and stabilise the ion. So more chlorine atoms (closer to the $\text{–COOH}$) give a stronger acid.

\section*{Esters}

An alcohol (or phenol) reacts with an acyl chloride at room temperature to give an \textbf{ester} 酯 and $\text{HCl}$. Examples are ethyl ethanoate (from ethanol) and phenyl benzoate (from phenol).

\section*{Acyl chlorides}

Acyl chlorides are made from carboxylic acids (with $\text{PCl}_3$, $\text{PCl}_5$ or $\text{SOCl}_2$). They are very reactive. At room temperature they react with:

\begin{center}
\adjustbox{max width=\linewidth}{%
\begin{tabular}{ll}
\toprule
\textbf{Reactant} & \textbf{Product (plus $\text{HCl}$)} \\
\midrule
water & the carboxylic acid \\
an alcohol & an ester \\
phenol & an ester \\
\textbf{ammonia} 氨 & an \textbf{amide} 酰胺 \\
a primary or secondary \textbf{amine} 胺 & an amide \\
\bottomrule
\end{tabular}}
\end{center}

\subsection*{The addition–elimination mechanism}

All these reactions follow an \textbf{addition–elimination} 加成消去 mechanism: a nucleophile first \textbf{adds} to the slightly positive carbonyl carbon, then $\text{HCl}$ is \textbf{eliminated}.

\subsection*{Ease of hydrolysis}

Compare how easily three chlorides react with water (\textbf{hydrolysis} 水解):

\[
\text{acyl chloride} \gg \text{alkyl chloride} \gg \text{aryl chloride}
\]

An acyl chloride reacts violently with cold water; an alkyl chloride reacts slowly; an aryl chloride (halogenoarene) does not react, because its C–Cl bond is strengthened by the ring.


\end{document}
